\section*{Chapter \ref{ch:b1:growth-accounting} --- Growth and Development Accounting}
% Numbers recomputed in code/bachelor-1/21-growth-accounting/exercises/solutions.py (tested).

\begin{solution}{exo:b1:growth-accounting:residual-ex}
$g_A = 4 - 0.35 \times 6 - 0.65 \times 1 = 4 - 2.1 - 0.65 = 1.25$ points, $31.25\%$ of growth.
\end{solution}

\begin{solution}{exo:b1:growth-accounting:read}
Germany $0.56/1.72 = 32.6\%$; Japan $0.3/2.33 = 12.9\%$.
\end{solution}

\begin{solution}{exo:b1:growth-accounting:steady}
Output per worker grows at $1.5\%$. Capital per worker grows at the same
rate, so capital deepening contributes $\alpha \times 1.5$ points and the residual
$(1 - \alpha) \times 1.5$: with $\alpha = 1/3$, $0.5$ and $1$ point.
\end{solution}

\begin{solution}{exo:b1:growth-accounting:china-ex}
Capital per worker: from $4.18\%$ to $8.53\%$ a year, a change of $4.35$
points; productivity from $-0.62\%$ to $3.18\%$, a change of $3.8$ points.
Capital per worker changed slightly more, and productivity changed sign.
(Output per worker, the sum the two terms explain, changed most: from
$2.1\%$ to $7.64\%$, $5.54$ points.)
\end{solution}

\begin{solution}{exo:b1:growth-accounting:levels-ex}
$y_{KH} = 0.2^{1/3} \times 0.6^{2/3} = 0.416$; productivity $0.25/0.416 = 0.601$; inputs account for
$\ln 0.416/\ln 0.25 = 63.3\%$ of the log gap.
\end{solution}

\begin{solution}{exo:b1:growth-accounting:two-firms}
$0.5l_1^{-0.5} = (100 - l_1)^{-0.5}$ gives $100 - l_1 = 4l_1$, so $l_1 = 20$ and $l_2 = 80$; the \omterm{def:b1:technology:marginal-product}{marginal
product} is $0.5 \times 20^{-0.5} = 0.1118$. Output rises from $3\sqrt{50} = 21.213$ to
$\sqrt{20} + 2\sqrt{80} = 22.361$, by $5.41\%$.
\end{solution}

\begin{solution}{exo:b1:growth-accounting:hours-ex}
Hours per worker fell in both countries, so the labour input grew more
slowly in hours than in persons: the labour contribution falls and the
residual, which is what is left, rises by the same amount. The fall was
larger in Germany.
\end{solution}

\begin{solution}{exo:b1:growth-accounting:quality}
The capital input grows faster than measured, so the contribution of
capital is understated and the residual overstated by $\alpha$ times the
unmeasured quality growth.
\end{solution}

\begin{solution}{exo:b1:growth-accounting:variance}
$100\%$: with $A$ equal everywhere, $\ln y = \ln A + \ln y_{KH}$ varies only through $\ln y_{KH}$. A
share of $29.6\%$ says that the dispersion of productivity, and its
correlation with inputs, account for most of the dispersion of income.
\end{solution}

\begin{solution}{exo:b1:growth-accounting:tornqvist}
With $\ln Y_t = \ln A_t + \alpha \ln K_t + (1 - \alpha)\ln(H_tL_t)$, differencing gives the decomposition
exactly. If $\alpha$ varies, write $\ln Y$ as a function of the log inputs; the
Törnqvist index is the trapezoid rule for the integral of the shares along
the path, exact for a translog function. Otherwise the second-order terms
cancel and the error is of the order of the cube of the yearly changes;
with one year's share in place of the average it would be of the order of
their square.
\end{solution}

\begin{solution}{exo:b1:growth-accounting:induced}
On the balanced path $\tilde k = K/(AL)$ is constant, so $K/L$ grows at $g$, and so does
$Y/L$. Accounting per worker: $g_y = \alpha g_k + g_{\text{res}}$ with $g_k = g$, so $g_{\text{res}} = (1 - \alpha)g$.
Yet without \omterm{def:b1:technology:technical-progress}{technical progress} capital per worker would not grow at all:
all growth is caused by technology.
\end{solution}

\begin{solution}{exo:b1:growth-accounting:refute}
Sustained growth of measured productivity in a planned economy, with
capital and labour inputs carefully measured --- for example output per
unit of input rising decade after decade at rates comparable with market
economies, or a return to capital that does not fall as capital piles up.
\end{solution}

\begin{solution}{pb:b1:growth-accounting:brenland}
\textbf{1.} $\ln(292390/197403)/25 = 1.571\%$ a year.

\textbf{2.} $\ln(4366/3760)/25 = 0.598\%$.

\textbf{3.} $1.571 - 0.598 = 0.974\%$ (with these rounded values $0.973$).

\textbf{4.} $K_{2000} = 592{,}209$; real investment $34200/77 \times 100 = 44{,}416$.

\textbf{5.} $K_{2025} = 953{,}616$; $g_K = \ln(953616/592209)/25 = 1.906\%$.

\textbf{6.} $953616/292390 = 3.261$, up from 3.

\textbf{7.} Capital $0.35 \times 1.906 = 0.667$; labour $0.65 \times 0.598 = 0.389$.

\textbf{8.} $1.571 - 0.667 - 0.389 = 0.516$ points (with these rounded values $0.515$), $32.8\%$ of growth.

\textbf{9.} Capital $0.572$, labour $0.418$.

\textbf{10.} $0.581$ points, $37\%$ of growth.

\textbf{11.} Lowering $\alpha$ moves weight from capital, which grew at $1.906\%$, to
labour, which grew at $0.598\%$; the inputs grow more slowly in all, and the
residual is larger.

\textbf{12.} Investment is the same, so a larger starting stock grows more
slowly in proportion: $g_K$ falls and the residual rises.

\textbf{13.} Errors in the deflator (real output), in the capital stock
(depreciation, starting value, quality) and in the labour input (hours,
skills).

\textbf{14.} Output 615,657, capital 757,491 (thousands of dalens) and 7,779
workers.

\textbf{15.} $11.39$ and $16.7$; the standard deviation of the logarithm is $0.149$.

\textbf{16.} Maximise $\sum A_i x_i^{0.85}$ with $x_i = k_i^{0.35}l_i^{0.65}$, subject to $\sum k_i = K$ and $\sum l_i = L$.
Equal marginal rates of substitution give every firm the same $k_i/l_i = K/L$,
so $x_i = s_i K^{0.35}L^{0.65}$ with $\sum s_i = 1$; equal \omterm{def:b1:technology:marginal-product}{marginal products} $0.85A_i s_i^{-0.15}$ give
$s_i \propto A_i^{1/0.15}$.

\textbf{17.} $Y^* = (\sum A_i^{1/0.15})^{0.15}(K^{0.35}L^{0.65})^{0.85} = 675{,}823$; the gain is $9.77\%$.

\textbf{18.} $13.81\%$: closer to constant returns, inputs move further.

\textbf{19.} $7.74\%$ of both inputs (firm 455).

\textbf{20.} Measurement error makes output, and so measured \omterm{def:b1:technology:marginal-product}{marginal
products}, more dispersed than the true ones: the computed gain includes
reallocations that would not raise true output.

\textbf{21.} $0.5^{1/3} \times 0.8^{2/3} = 0.684$.

\textbf{22.} $0.6/0.684 = 0.877$.

\textbf{23.} $\ln 0.684/\ln 0.6 = 74.4\%$.

\textbf{24.} Across the world in 2019, inputs account for only $29.6\%$ of the
variance of $\ln y$ (a different measure, over many countries); Kestria's gap
with Brenland is mostly one of inputs.

\textbf{25.} Closing the capital gap raises income by the factor
$2^{1/3} = 1.26$; closing the productivity gap by $1/0.877 = 1.14$. Capital first,
though capital takes years to build.

\textbf{26.} Equalising \omterm{def:b1:technology:marginal-product}{marginal products} in Brenland's census would raise
output by $9.8\%$ (with \omterm{def:b1:technology:returns-to-scale}{returns to scale} of $0.85$; $13.8\%$ with $0.9$).
\end{solution}
